\documentclass[10pt,a4paper]{article}

% --- Packages ---
\usepackage[utf8]{inputenc}
\usepackage[T1]{fontenc}
\usepackage{lmodern}
\usepackage[margin=2.5cm]{geometry}
\usepackage{amsmath,amssymb,amsfonts,bm}
\usepackage{booktabs}
\usepackage{array}
\usepackage{graphicx}
\usepackage{cite}
\usepackage{microtype}
\usepackage{authblk}
\usepackage{abstract}
\usepackage[hidelinks,unicode,pdfcreator={LaTeX}]{hyperref}

\hypersetup{
  pdftitle={S4D-TAM: Semantic 4D Token Attention Map for Autonomous Navigation of Unmanned Aerial Vehicles in GNSS-Degraded Environments},
  pdfauthor={Mateusz Pomianek},
  pdfkeywords={unmanned aerial vehicles, semantic SLAM, 4D tokens, spatiotemporal transformer, autonomous navigation, GNSS degradation, world model, occupancy forecasting}
}

% --- Document Information ---
\title{\textbf{S4D-TAM: Semantic 4D Token Attention Map for Autonomous Navigation of Unmanned Aerial Vehicles in GNSS-Degraded Environments}}

\author{Mateusz Pomianek}
\affil{Department of Computer and Control Engineering, Rzeszów University of Technology, Poland\\
\texttt{m.pomianek@prz.edu.pl}}

\date{\today}

\begin{document}

\maketitle

\begin{abstract}
We propose S4D-TAM (Semantic 4D Token Attention Map), a novel navigation system architecture for unmanned aerial vehicles (UAVs) operating in GNSS-degraded environments. Unlike classical methods that repeatedly analyze raw image streams, the proposed system transforms sensor data (RGB, thermal, IMU, GNSS) into a persistent map of semantic 4D tokens. Each token represents a spatial fragment or object instance and contains its geometry, semantic embedding, dynamics, observation history, and calibrated uncertainty. A hierarchical spatiotemporal transformer updates the world model, aligns it with reference maps, and forecasts occupancy in partially observable regions. The predictive map is then used by a planner that accounts for risk, energy cost, and information value of future observations. The proposed representation aims to reduce image processing redundancy, increase robustness to GNSS degradation, and enable decision-making based on persistent spatial memory rather than solely on the current sensor field of view. The article presents the system concept, literature review, mathematical description, an executable auditable reference implementation, and a two-level experimental methodology separating external system comparison from internal mechanism ablation. Confirmatory experimental results are not yet reported.
\end{abstract}

\textbf{Keywords:} unmanned aerial vehicles, semantic SLAM, 4D tokens, spatiotemporal transformer, autonomous navigation, GNSS degradation, world model, occupancy forecasting

\vspace{1em}
\hrule
\vspace{1em}

% --- Content Sections ---
\section{Introduction}
\label{sec:introduction}

Autonomous navigation of unmanned aerial vehicles (UAVs) in environments with limited GNSS signal availability is one of the central challenges of modern mobile robotics. Satellite signal degradation can occur in dense urban environments, under heavy forest canopy, in tunnels and canyons, and under deliberate jamming or spoofing conditions~\cite{yue2025airground, radwan2024uav, delmerico2018benchmark, qin2018vinsmono}. In such scenarios, a UAV must rely on alternative information sources: RGB and thermal cameras, inertial measurement units (IMU), altimeters, LiDARs, and reference maps~\cite{nieuwenhuisen2015autonomous, nieuwenhuisen2016autonomous, yue2024semantic, hu2025review}.

Contemporary autonomous navigation systems perform pose estimation, map construction, semantic segmentation, object recognition, occupancy forecasting, and trajectory planning as separate, loosely coupled modules~\cite{feng2025survey, placed2022survey}. This approach leads to computational redundancy -- the same images are repeatedly analyzed by different modules, and semantic and geometric information is not efficiently shared between perception and planning stages~\cite{rosen2024scene, ding2025survey}.

In this work, we propose S4D-TAM (Semantic 4D Token Attention Map), a novel navigation system architecture in which the fundamental environment representation is a persistent, semantic, continuously updated map of 4D tokens derived from the SLAM point cloud. RGB and thermal camera data are used primarily to update the geometry and semantic meaning of points and objects stored in the map. Planning, risk assessment, and decision-making modules do not repeatedly analyze the same images -- they operate on a compressed spatial representation containing geometry, semantics, dynamics, observation history, and uncertainty~\cite{pomianek2026s4dtam}.

The key scientific novelty of the proposed system is the introduction of a persistent memory map composed of semantic 4D tokens, which serves as a unified representation for localization, environment understanding, future scene state prediction, risk estimation, and UAV trajectory planning. The system remembers not only where an object was located, but also when and how many times it was observed, from which sensors the information originated, how reliable its position and recognition are, whether the object is static, dynamic, or periodically displaced, what spatial relations hold between it and other terrain elements, what invisible areas may be hidden behind it, and what significance the object has for flight safety and planning~\cite{pomianek2026s4dtam}.

The article combines a conceptual system architecture with an executable research reference and a reproducible benchmark framework. The current repository implements token state and lifecycle management, proposal and association, multimodal encoder interfaces, attention-related processing, uncertainty calibration utilities, reference-map and topological support, causal occupancy and motion forecasting, deterministic risk-aware planning, telemetry, and the integrated \texttt{S4DTAMReference} execution path. The final learned hierarchical model, full public-dataset reproduction, and confirmatory results remain under development. Section~\ref{sec:experiments} therefore specifies a prospective two-level validation protocol rather than reporting final performance results.

\section{Literature Review}
\label{sec:literature_review}

\subsection{Semantic SLAM for UAVs}
\label{subsec:semantic_slam}

Semantic SLAM for aerial platforms must be lightweight, robust to GPS signal loss, and capable of producing useful semantic maps for planning and multi-robot coordination. Liu et al. proposed SlideSLAM -- a sparse, lightweight, decentralized metric-semantic SLAM for multi-robot navigation that uses an object-based representation and enables efficient map merging in heterogeneous aerial-ground teams~\cite{liu2024slideslam}. Tao et al. integrated active exploration with metric-semantic SLAM for resource-constrained aerial robots, explicitly balancing exploration and localization uncertainty (including semantic loop closures), which informs S4D-TAM strategies for uncertainty-aware mapping in GNSS-denied environments~\cite{tao2024active}. Liu et al. presented a real-time semantic SLAM that detects and models tree trunks and planes from LiDAR to constrain poses and enable drift compensation during long-range autonomous flights under forest canopy, directly relevant to drift mitigation and semantic constraints in S4D-TAM trajectory planning~\cite{liu2022largescale}.

Tian et al. proposed SAME -- a ground-air collaboration system that fuses semantic Octree units from a UAV and a UGV into a coherent global map supporting collaborative exploration, offering a semantic octree architecture and fusion approach useful in S4D-TAM for cross-platform mapping and multi-perspective scene reconstruction~\cite{tian2024same}. Xiang and Su developed a LiDAR-based semantic mapping pipeline that creates hierarchical semantic maps together with occupancy grids for navigation and task adaptation, illustrating practical semantic mapping designs and hierarchical semantics that S4D-TAM can exploit for task-oriented scene understanding~\cite{xiang2021customizable}.

Li et al. introduced Hi-SLAM -- a hierarchical categorical representation within 3D Gaussian Splatting for compact encoding of multiple semantic classes and improved semantic mapping and tracking, directly applicable to S4D-TAM requirements for scalable semantic scene representations~\cite{li2024hislam, keetha2024splatam}. More recent semantic SLAM works employing 3D Gaussian Splatting include SGS-SLAM~\cite{hu2024sgsslam}, GS4/GaussianSLAM~\cite{yugay2023gaussianslam}, SemGauss-SLAM~\cite{zou2024semgaussslam}, Photo-SLAM~\cite{huang2024photoslam}, DynaGS-SLAM~\cite{wu2024dynags}, GSFF-SLAM~\cite{pan2024gsffslam}, OVO-SLAM~\cite{haghighi2024ovoslam}, MSGS-SLAM~\cite{li2024msgsslam}, NIS-SLAM~\cite{li2024nisslam}, and NEDS-SLAM~\cite{dai2024nedsslam}, demonstrating growing interest in real-time dense semantic representations for SLAM.

\subsection{Point Cloud Transformers}
\label{subsec:point_cloud_transformers}

Point cloud transformer architectures and related models provide expressive spatiotemporal encoding of 3D geometry and dynamics, useful for the point-cloud-based perception and temporal prediction modules of S4D-TAM. PCPNet introduces a transformer network for future LiDAR point cloud prediction with a semantic auxiliary training objective for semantically consistent forecasts, directly applicable to S4D-TAM for point-level prediction and semantic alignment~\cite{pcpnet2024}. Silva et al. proposed S2TPVFormer, which extends TPV embeddings with hybrid cross-view temporal attention for temporally consistent semantic 3D occupancy, offering a spatiotemporal transformer design that S4D-TAM can adopt for temporally consistent semantic occupancy mapping~\cite{silva2024s2tpvformer}.

Chen et al. presented AdaOcc, which employs adaptive forward-view transformations and flow modeling (with strong Swin camera backbone models) to improve voxel occupancy and flow prediction, providing practical camera-to-3D transformer strategies applicable in vision-oriented S4D-TAM modules~\cite{chen2024adaocc}. Work on point cloud forecasting as a proxy for 4D occupancy prediction redefines the problem toward sensor-agnostic spatiotemporal (4D) occupancy forecasting, directing S4D-TAM toward world-oriented 4D occupancy targets for robust temporal modeling~\cite{yu2025driveoccworld}.

Further relevant works include Point-BERT~\cite{yu2022pointbert}, which introduces masked point modeling for pre-training point cloud transformers; SAT3D / Point Attention Network~\cite{ji2020point} for semantic segmentation of 3D scenes using slot attention; and YOGO~\cite{zhang2021you} and Efficient Point Transformer~\cite{he2024efficient, guo2021pct}, demonstrating efficient point cloud processing with token representation and relation inference modules relevant to S4D-TAM token compression.

\subsection{4D Occupancy Forecasting}
\label{subsec:occupancy_forecasting}

Accurate 4D occupancy forecasting provides the spatiotemporal scene model required by S4D-TAM for risk assessment and trajectory planning. Yu et al. proposed Drive-OccWorld, which adapts a vision-centric 4D world model with memory normalization and action-conditioned decoding for occupancy and flow prediction, and integrates forecasting with end-to-end planning, providing an approach that S4D-TAM can leverage to couple occupancy prediction with trajectory selection~\cite{yu2025driveoccworld}. Zheng et al. presented OccWorld -- a spatiotemporal generative transformer architecture for 4D occupancy prediction, demonstrating effectiveness on driving datasets and offering S4D-TAM a transformer backbone for long-horizon volumetric prediction~\cite{zheng2024occworld}.

Ma et al. provided Cam4DOcc -- a standardized camera-only benchmark and baselines for current and future occupancy estimation, providing S4D-TAM with evaluation protocols and camera-only baseline designs for 4D occupancy forecasting~\cite{ma2023cam4docc}. OccSora uses a 4D scene tokenizer and a diffusion transformer for generative long-horizon 4D occupancy simulation conditioned on trajectories, offering S4D-TAM a generative world simulator approach for scenario generation and planning evaluation~\cite{wang2024occsora}. Liao et al. proposed I2-World with hierarchical intra/inter tokenizers and an encoder-decoder design for efficient compression and generation of temporally consistent 4D occupancy, presenting a tokenization strategy that S4D-TAM can apply for compact world representations and fast inference~\cite{liao2025i2world}.

Zheng et al. presented GaussianAD, which represents scenes through semantic 3D Gaussians refined from images, predicts 3D flows for those Gaussians, and integrates forecasting with planning -- this sparse semantic Gaussian approach is directly relevant to S4D-TAM's lightweight scene abstractions and prediction-based planning~\cite{zheng2024gaussianad}.

\subsection{World Models and Planning}
\label{subsec:world_models}

World models unify perception, prediction, and planning; for S4D-TAM, they provide a mechanism for simulating the future, evaluating trajectories, and incorporating action-conditioned forecasts. Chen et al. proposed DriveWorld -- a Memory State-Space Model with dynamic memory and static scene propagation for spatiotemporal pre-training, demonstrating downstream gains in detection, mapping, tracking, forecasting, and planning -- relevant to S4D-TAM pre-training strategies and multi-task world modeling~\cite{chen2024driveworld}.

Zhao et al. used a contrastive world model objective to learn visual representations for drone navigation and visuomotor policy learning, illustrating how world model pre-training can improve navigation policies, which S4D-TAM can leverage for improved visual generalization~\cite{zhao2023learning}. Mei et al. proposed IR-WM, which focuses on predicting residual changes relative to a temporal BEV prior and couples forecasting with planning, offering S4D-TAM a compact residual prediction approach that reduces redundant static background modeling and improves planning integration~\cite{mei2025irwm}.

Cai et al. presented NEUSIS -- a compositional neuro-symbolic framework combining neuro-symbolic perception, a probabilistic world model with Bayesian filtering, and hierarchical search planning, exemplifying an interpretable and uncertainty-aware world modeling and planning workflow useful for S4D-TAM mission-level decision-making~\cite{cai2024neusis}. Liu et al. proposed HPP, which integrates occupancy-conditioned prediction with game-theoretic motion forecasting and planning, demonstrating methods for aligning prediction patterns with planning objectives that S4D-TAM can adopt for improved trajectory selection under uncertainty and interaction~\cite{liu2025hybrid}.

Feng et al. provided a comprehensive survey of world models for autonomous driving, covering 4D occupancy forecasting, generative simulation, and planning integration, providing a conceptual foundation and taxonomy that can guide S4D-TAM design choices in world model selection, pre-training, and evaluation~\cite{feng2025survey}.

\subsection{UAV Navigation in GNSS-Denied Environments}
\label{subsec:gnss_denied}

Autonomous UAV navigation in GNSS-degraded environments requires multi-sensor fusion, robust SLAM algorithms, and uncertainty-aware planning strategies. Yue et al. proposed an air-ground cooperative UAV autonomous navigation system for GPS-denied environments that combines visual odometry, IMU, and communication with a ground robot for robust localization~\cite{yue2025airground}. Radwan et al. presented UAV-assisted visual SLAM generating reconstructed 3D scene graphs in GPS-denied environments, demonstrating scene graph representations for scene understanding and planning~\cite{radwan2024uav}.

Nieuwenhuisen et al. developed an autonomous navigation system for micro aerial vehicles in complex GNSS-denied environments, using LiDAR, IMU, and obstacle-aware planning~\cite{nieuwenhuisen2015autonomous, nieuwenhuisen2016autonomous}. Yue et al. proposed semantically guided visual autonomous navigation for UAVs that employs semantic segmentation to identify navigable areas and obstacles for improved path planning~\cite{yue2024semantic}. Hu et al. provided a review of vision-based UAV localization and optimization-based path planning in GPS-denied environments, synthesizing SLAM, object detection, and trajectory planning methods~\cite{hu2025review}.

Further relevant works include UAV navigation systems using multi-sensor fusion (IMU, camera, LiDAR, radar)~\cite{delmerico2018benchmark, qin2018vinsmono, shan2020liosam, zhang2014loam, xu2022fastlio2, faessler2015automatic}, KAZE feature-based SLAM~\cite{alcantarilla2013kaze}, map-aided navigation and dead reckoning~\cite{brossard2020aiimu}, and platforms for drone racing and agile navigation~\cite{delmerico2019uzhfpv}. Castano presented autonomous UAV navigation in GPS-denied environments using LiDAR point clouds for localization and mapping~\cite{castano2023autonomous}.

\section{Research Gap and Motivation}
\label{sec:research_gap}

Contemporary solutions perform pose estimation, map construction, semantic segmentation, object recognition, occupancy prediction, trajectory planning, and risk assessment as separate, independent modules~\cite{feng2025survey, placed2022survey, ding2025survey}. Systems such as LVI-SAM integrate LiDAR, camera, and IMU in a factor graph, while Kimera constructs metric-semantic 3D models. Neither, however, constitutes a complete, learned world model that exploits persistent map tokens for predicting unseen hazards and guiding sensor attention~\cite{rosen2024scene}.

OpenScene enables assigning language-image embedding space features to point cloud points. Open3DSG extends such approaches with inter-object relationships, while Octree-Graph combines semantics, occupancy, and graph structure. These solutions, however, primarily address scene understanding rather than closed-loop UAV control operating under limited computational resources and partial observability~\cite{tian2024same, xiang2021customizable}.

Research on 4D occupancy forecasting has demonstrated that predicting the future state of space can serve as an effective planning foundation~\cite{yu2025driveoccworld, mei2025irwm, chen2024driveworld, zheng2024gaussianad}. Nevertheless, automotive applications, dense voxel representations, and full sensor sequence processing still dominate. The proposed S4D-TAM system fills the gap between semantic SLAM, 4D world models, and risk-aware planning for UAVs in GNSS-degraded environments.

The key motivation is to transition from a classical geometric map to a predictive spatial memory of the UAV -- one that not only records where objects are, but also predicts where they may appear, which areas are unobservable, and what risk they pose to flight safety. This approach can be termed a \emph{SLAM-native world model}: a world model whose primary representation space is the SLAM map rather than the raw image stream~\cite{pomianek2026s4dtam}.

\section{S4D-TAM System Concept}
\label{sec:system_concept}

\subsection{Definition of the 4D Token}
\label{subsec:token_definition}

A standard 3D token represents a spatial fragment or object at a given position:
\begin{equation}
  t_i = [x_i, y_i, z_i]
\end{equation}
However, it cannot capture environmental changes or observation history. In an autonomous system, it is essential to distinguish a permanent structure (building, tree, or terrain) from dynamic or movable objects (persons, vehicles) that may change their position over time~\cite{pomianek2026s4dtam}.

A 4D token incorporates the temporal dimension alongside geometric, semantic, dynamic, and uncertainty attributes:
\begin{equation}
  T_i(t) = [\mathbf{p}_i, \mathbf{\Sigma}_i, \mathbf{g}_i, \mathbf{s}_i, \mathbf{v}_i, \mathbf{h}_i, u_i, r_i]
\end{equation}
where the individual elements denote:
\begin{enumerate}
  \item $\mathbf{p}_i \in \mathbb{R}^3$ -- position in 3D space,
  \item $\mathbf{\Sigma}_i \in \mathbb{R}^{3 \times 3}$ -- position covariance matrix (geometric uncertainty),
  \item $\mathbf{g}_i \in \mathbb{R}^{d_g}$ -- geometric descriptor (surface normals, curvature, local geometric features),
  \item $\mathbf{s}_i \in \mathbb{R}^{d_s}$ -- semantic embedding (features from a deep neural network or open-vocabulary foundation model),
  \item $\mathbf{v}_i \in \mathbb{R}^3$ -- estimated velocity and motion vector,
  \item $\mathbf{h}_i \in \mathbb{R}^{d_h}$ -- observation history (observation count, timestamps, sensor modality sources),
  \item $u_i \in [0, 1]$ -- overall token uncertainty (scalar reliability measure),
  \item $r_i \in [0, 1]$ -- local risk indicator (hazard probability).
\end{enumerate}

The fourth dimension does not merely append time as an extra coordinate; the token encodes persistence, observation chronology, rate of change, and the forecasted future state of the object~\cite{pomianek2026s4dtam, zheng2024gaussianad, wang2024occsora, ma2023cam4docc, chen2024riskoccupancy}.

\subsection{System Architecture}
\label{subsec:architecture}

The S4D-TAM system is organized into nine functional layers:
\begin{itemize}
  \item \textbf{Layer 1: Sensor Fusion} -- Integration of data from RGB and thermal cameras, IMU, GNSS, altimeter, and optionally LiDAR, radar, and prior reference maps. Each measurement is associated with a timestamp, coordinate frame transformation, and sensor reliability measure~\cite{pomianek2026s4dtam}.
  \item \textbf{Layer 2: SLAM Front-End} -- A lightweight visual-inertial front-end performing feature tracking, depth estimation, landmark detection, and feature vector extraction. The geometric SLAM constructs a local point cloud and tracks the UAV trajectory. The factor graph incorporates visual odometry, inertial, GNSS, altimeter, semantic loop-closure, and reference-map alignment factors~\cite{liu2024slideslam, tao2024active, liu2022largescale}.
  \item \textbf{Layer 3: Tokenization and Hierarchical Representation} -- The SLAM point cloud is partitioned into superpoints, object segments, adaptive octree cells, local surface patches, and dynamic instances. Each segment is encoded into one or more tokens. Large regular surfaces are stored at lower resolution, while edges, obstacles, and dynamic objects are retained at higher resolution~\cite{li2024hislam, tian2024same, xiang2021customizable, zhang2021you, he2024efficient}.
  \item \textbf{Layer 4: Spatiotemporal Transformer} -- A hierarchical attention mechanism combining local scope (tokens within a spatial radius, predicted flight corridor, sensor field of view) and global scope (landmarks, major terrain obstacles, dynamic objects, high-uncertainty regions, elevated-risk locations)~\cite{yu2025driveoccworld, silva2024s2tpvformer, liao2025i2world, pcpnet2024}.
  \item \textbf{Layer 5: Object Memory and Semantics} -- Upon detecting a landmark or object, an instance token is instantiated storing class probability distributions, observation history, open-vocabulary embeddings, and spatial relations with neighboring objects~\cite{li2024hislam, cai2024neusis, haghighi2024ovoslam}.
  \item \textbf{Layer 6: Reference Map Alignment} -- Maintenance and alignment of the online map $M_t$ with an a priori reference map $M_{\mathrm{ref}}$ (elevation maps, prior mission clouds) via a three-stage procedure: global topological descriptor matching, landmark token matching, and local geometric optimization~\cite{liu2022largescale, yue2025airground, radwan2024uav}.
  \item \textbf{Layer 7: Occupancy and Risk Forecasting} -- Prediction of future space occupancy and hazard probability in partially observable areas (inferring occluded regions behind building edges, terrain trails leading to unseen objects, and dynamic object trajectories)~\cite{yu2025driveoccworld, mei2025irwm, liao2025i2world, zheng2024gaussianad, liu2025hybrid}.
  \item \textbf{Layer 8: Trajectory Planning} -- A risk-aware planner optimizing an objective function balancing path length, energy consumption, collision probability, predicted environmental risk, and uncertainty penalty~\cite{cai2024neusis, yue2025airground, tordesillas2022faster}.
  \item \textbf{Layer 9: Active Perception} -- Dynamic control of sensor orientation, processing frequency, and sensory modality based on expected information gain (uncertainty reduction) relative to operational cost~\cite{pomianek2026s4dtam, tao2024active, tian2024same}.
\end{itemize}

\subsection{Current Reference Implementation and Target Research Modules}
\label{subsec:implementation_status}

The S4D-TAM concept is accompanied by an executable Python research reference. The implementation is modular so that individual mechanisms can be tested, replaced, and ablated without altering evaluator contracts.

Implemented reference components include persistent token state and lifecycle management (\texttt{token.py}, \texttt{memory.py}), token proposal and data association (\texttt{proposal.py}, \texttt{association.py}), modality-oriented encoder interfaces and masked fusion (\texttt{encoders/}), attention-related processing (\texttt{attention.py}), calibration and uncertainty utilities (\texttt{calibration.py}), reference-map and topological support (\texttt{reference\_map.py}, \texttt{topology.py}), causal probabilistic occupancy and motion forecasting (\texttt{forecasting.py}), deterministic risk-, energy-, time-, progress-, and information-aware planning (\texttt{planner.py}), structured event telemetry (\texttt{telemetry.py}), and the integrated \texttt{S4DTAMReference} pipeline (\texttt{pipeline.py}).

The reference implementation favors strict causality, deterministic replay, explicit uncertainty, and auditability over raw throughput. During early numerical validation, observed differences can be traced to defined mechanisms rather than stochastic non-determinism.

Formulations in Sections~\ref{subsec:bayesian_inference}--\ref{subsec:resource_management} specify target research extensions alongside the reference implementation, including learned hierarchical transformer backends, multimodal DBN/particle filters, MPPI stochastic trajectory sampling, adaptive octree Level-of-Detail (LOD) memory management, and real-time EDF scheduling.

\section{Mathematical Description}
\label{sec:math_description}

\subsection{Sensor Layer and Data Fusion}
\label{subsec:sensor_fusion}

The system employs a multi-modal sensor suite $\mathcal{S} = \{s_1, s_2, \dots, s_n\}$, where each sensor $s_i$ provides a measurement $\mathbf{z}_i(t)$ at time $t$ characterized by covariance $\mathbf{\Sigma}_i(t)$. Sensor data are tightly fused within a factor graph formulation:
\begin{equation}
  F_{\text{total}} = F_{\text{VO}} + F_{\text{IMU}} + F_{\text{GNSS}} + F_{\text{alt}} + F_{\text{sem}} + F_{\text{ref}}
\end{equation}
where the individual factors correspond to:
\begin{itemize}
  \item $F_{\text{VO}}$ -- visual odometry factor,
  \item $F_{\text{IMU}}$ -- inertial measurement preintegration factor,
  \item $F_{\text{GNSS}}$ -- GNSS position factor (with adaptive covariance under signal degradation),
  \item $F_{\text{alt}}$ -- barometric / LiDAR altimeter factor,
  \item $F_{\text{sem}}$ -- semantic loop-closure factor,
  \item $F_{\text{ref}}$ -- reference map alignment factor.
\end{itemize}
Each factor is formulated as a Mahalanobis distance penalty:
\begin{equation}
  F_i = \sum_k \| \mathbf{e}_{ik} \|_{\mathbf{\Sigma}_{ik}}^2
\end{equation}
where $\mathbf{e}_{ik}$ is the residual error and $\mathbf{\Sigma}_{ik}$ is the corresponding measurement covariance matrix~\cite{liu2024slideslam, tao2024active, liu2022largescale}.

\subsection{SLAM Front-End}
\label{subsec:slam_frontend}

The SLAM front-end jointly estimates the UAV trajectory $\mathbf{X} = \{\mathbf{x}_1, \mathbf{x}_2, \dots, \mathbf{x}_T\}$ and the spatial point map $\mathcal{M} = \{\mathbf{m}_1, \mathbf{m}_2, \dots, \mathbf{m}_N\}$ by minimizing the robustified cost:
\begin{equation}
  \mathbf{X}^*, \mathcal{M}^* = \arg\min_{\mathbf{X}, \mathcal{M}} \sum_{i,j} \rho\left( \| \mathbf{z}_{ij} - h(\mathbf{x}_i, \mathbf{m}_j) \|_{\mathbf{\Sigma}_{ij}}^2 \right)
\end{equation}
where $h(\mathbf{x}_i, \mathbf{m}_j)$ denotes the projective measurement function of 3D landmark $\mathbf{m}_j$ from pose $\mathbf{x}_i$, and $\rho(\cdot)$ is a robust M-estimator kernel (e.g., Huber loss)~\cite{liu2024slideslam, liu2022largescale}.

Metric scale is resolved and stabilized via:
\begin{equation}
  s^* = \arg\min_s \sum_k w_k \| s \cdot d_k^{\text{VO}} - d_k^{\text{ref}} \|^2
\end{equation}
where $d_k^{\text{VO}}$ is the distance measured by visual odometry, $d_k^{\text{ref}}$ is the metric reference baseline obtained from IMU double-integration, altimeter, GNSS, known object dimensions, or digital elevation models, and $w_k$ denotes the reliability weight~\cite{pomianek2026s4dtam, liu2022largescale}.

\subsection{Point Cloud Tokenization}
\label{subsec:tokenization}

The continuous point cloud $\mathcal{M}$ is partitioned into $K$ spatial tokens $\{T_1, T_2, \dots, T_K\}$ via hierarchical segmentation:
\begin{equation}
  \mathcal{M} = \bigcup_{i=1}^K \mathcal{P}_i
\end{equation}
where $\mathcal{P}_i$ is the point cluster associated with token $T_i$. Partitioning strategies include:
\begin{enumerate}
  \item Superpoints -- grouping points of homogeneous geometry and semantics~\cite{zhang2021you, he2024efficient},
  \item Object segments -- instance bounding boxes and 3D masks~\cite{li2024hislam, haghighi2024ovoslam},
  \item Adaptive octrees -- hierarchical volumetric spatial decomposition~\cite{tian2024same, xiang2021customizable},
  \item Local surface patches -- planar surfaces, edges, and corners~\cite{liu2022largescale},
  \item Dynamic instances -- detected moving entities~\cite{zheng2024gaussianad, wang2024occsora}.
\end{enumerate}

For each token $T_i$, canonical spatial and semantic attributes are extracted:
\begin{align}
  \mathbf{p}_i &= \frac{1}{|\mathcal{P}_i|} \sum_{\mathbf{m} \in \mathcal{P}_i} \mathbf{m} \\
  \mathbf{\Sigma}_i &= \frac{1}{|\mathcal{P}_i|} \sum_{\mathbf{m} \in \mathcal{P}_i} (\mathbf{m} - \mathbf{p}_i)(\mathbf{m} - \mathbf{p}_i)^\top \\
  \mathbf{g}_i &= \mathrm{GeomEncoder}(\mathcal{P}_i) \\
  \mathbf{s}_i &= \mathrm{SemEncoder}(\mathcal{I}_i, \mathcal{P}_i)
\end{align}
where $\mathcal{I}_i$ is the set of image views observing point set $\mathcal{P}_i$~\cite{li2024hislam, pcpnet2024, he2024efficient}.

\subsection{Spatiotemporal Attention Mechanism}
\label{subsec:attention}

The hierarchical spatiotemporal transformer refines token embeddings via self- and cross-attention:
\begin{equation}
  T_i' = T_i + \mathrm{Attention}(\mathbf{Q}_i, \mathbf{K}, \mathbf{V})
\end{equation}
with linear projections $\mathbf{Q}_i = \mathbf{W}_Q T_i$, $\mathbf{K} = \mathbf{W}_K T$, and $\mathbf{V} = \mathbf{W}_V T$.

Local attention restricts query keys to the spatial neighborhood:
\begin{equation}
  \mathcal{N}_i^{\text{local}} = \{j : \|\mathbf{p}_i - \mathbf{p}_j\| < r_{\text{local}}\}
\end{equation}
while global attention incorporates key landmarks, dynamic entities, and critical safety hazards:
\begin{equation}
  \mathcal{N}_i^{\text{global}} = \{j : \mathrm{landmark}(j) \lor \mathrm{dynamic}(j) \lor u_j > \theta_u \lor r_j > \theta_r\}
\end{equation}

The unified attention weights incorporate spatial, temporal, and uncertainty gates:
\begin{equation}
  \alpha_{ij} = \frac{\exp\left( \frac{\mathbf{Q}_i^\top \mathbf{K}_j}{\sqrt{d_k}} \cdot \varphi(\mathbf{p}_i, \mathbf{p}_j, t_i, t_j) \cdot \psi(u_i, u_j) \right)}{\sum_{k \in \mathcal{N}_i} \exp\left( \frac{\mathbf{Q}_i^\top \mathbf{K}_k}{\sqrt{d_k}} \cdot \varphi(\mathbf{p}_i, \mathbf{p}_k, t_i, t_k) \cdot \psi(u_i, u_k) \right)}
\end{equation}
where $\varphi(\cdot)$ embeds relative spatial displacement and time intervals, while $\psi(\cdot)$ modulates attention based on observation confidence~\cite{yu2025driveoccworld, silva2024s2tpvformer, liao2025i2world, pcpnet2024}.

\subsection{Object Memory and Semantics}
\label{subsec:object_memory}

Upon detecting a recognized object instance, an instance token is instantiated:
\begin{equation}
  O_k = [\mathbf{p}_k, \mathbf{\Sigma}_k, \mathbf{b}_k, \mathbf{s}_k, P(c_k \mid \mathbf{z}_{1:t}), \mathbf{h}_k]
\end{equation}
where $\mathbf{b}_k$ is the 3D bounding box, and $P(c_k \mid \mathbf{z}_{1:t})$ is the class posterior probability over categorical classes $c_k$.

The semantic class distribution is updated recursively via Bayesian filtering:
\begin{equation}
  P(c_k \mid \mathbf{z}_{1:t}) = \frac{P(\mathbf{z}_t \mid c_k) P(c_k \mid \mathbf{z}_{1:t-1})}{\sum_{c'} P(\mathbf{z}_t \mid c') P(c' \mid \mathbf{z}_{1:t-1})}
\end{equation}
Open-vocabulary embeddings provide zero-shot classification capabilities for novel obstacle classes without closed-set retraining~\cite{li2024hislam, cai2024neusis, haghighi2024ovoslam}.

\subsection{Reference Map Alignment}
\label{subsec:refmap_alignment}

The system co-registers the active local map $M_t$ with an a priori reference map $M_{\text{ref}}$ via a three-tier optimization scheme:
\begin{enumerate}
  \item \textbf{Global Alignment} -- Coarse transform estimation via global descriptor matching:
  \begin{equation}
    T_{\text{global}} = \arg\min_T \sum_i \| \mathbf{d}_i^t - \mathbf{d}_{\sigma(i)}^{\text{ref}} \|^2
  \end{equation}
  where $\mathbf{d}_i^t$ and $\mathbf{d}_{\sigma(i)}^{\text{ref}}$ denote topological and elevation descriptors~\cite{liu2022largescale, yue2025airground, radwan2024uav}.
  \item \textbf{Landmark Matching} -- Transform refinement using prominent semantic landmarks $\mathcal{L}$:
  \begin{equation}
    T_{\text{landmark}} = \arg\min_T \sum_{j \in \mathcal{L}} \| T \mathbf{p}_j^t - \mathbf{p}_{\sigma(j)}^{\text{ref}} \|_{\mathbf{\Sigma}_j}^2
  \end{equation}
  \item \textbf{Local Geometric Optimization} -- Fine registration via semantic ICP:
  \begin{equation}
    T_{\text{final}} = \arg\min_T \sum_{i,j} w_{ij} \| T \mathbf{p}_i^t - \mathbf{p}_j^{\text{ref}} \|^2
  \end{equation}
  where $w_{ij}$ enforces semantic compatibility between matched patches~\cite{liu2022largescale, li2024hislam}.
\end{enumerate}

Map discrepancies provide direct diagnostic cues: structural changes, newly appeared obstacles, or localized GNSS measurement offsets~\cite{pomianek2026s4dtam, yue2025airground}.

\subsection{Occupancy and Risk Forecasting}
\label{subsec:occupancy_forecasting_math}

Future volumetric occupancy and operational risk are forecasted forward over horizon $\Delta t$:
\begin{equation}
  \mathcal{O}(t+\Delta t) = f_{\text{pred}}(T_{1:K}(t), \mathbf{a}(t))
\end{equation}
where $\mathcal{O}(t+\Delta t)$ represents forecasted occupancy grids/voxels, $T_{1:K}(t)$ is the current token set, and $\mathbf{a}(t)$ is the planned UAV action~\cite{yu2025driveoccworld, mei2025irwm, liao2025i2world, zheng2024gaussianad}.

The spatial risk map is defined as:
\begin{equation}
  R(\mathbf{p}, t+\Delta t) = P(H_{\mathbf{p}, t+\Delta t} \mid M_{1:t}, M_{\text{ref}}, C_t)
\end{equation}
where $H_{\mathbf{p}, t+\Delta t}$ denotes a collision hazard at position $\mathbf{p}$ and $C_t$ represents environmental context~\cite{pomianek2026s4dtam, liu2025hybrid}.

Occluded region risk is inferred from semantic contextual priors:
\begin{equation}
  R_{\text{occluded}}(\mathbf{p}) = \sum_{c \in \mathcal{C}} P(c \mid T_{\text{visible}}) \cdot P(H \mid c, \mathrm{geometry}(\mathbf{p}))
\end{equation}
where $\mathcal{C}$ encapsulates contextual categories (e.g., roads, intersections, building perimeters)~\cite{pomianek2026s4dtam, cai2024neusis, liu2025hybrid}.

\subsection{Trajectory Planning}
\label{subsec:trajectory_planning}

The trajectory planner computes an optimal control trajectory $\xi(t)$ minimizing a composite functional:
\begin{equation}
  J(\xi) = \int_0^T \left[ L(\xi(t)) + E(\xi(t), \dot{\xi}(t)) + \lambda_H P_H(\xi(t)) + \lambda_R R(\xi(t), t) + \lambda_U U(\xi(t)) \right] dt
\end{equation}
where:
\begin{itemize}
  \item $L(\xi(t))$ -- trajectory path length,
  \item $E(\xi(t), \dot{\xi}(t))$ -- aerodynamic and propulsion energy expenditure,
  \item $P_H(\xi(t)) = 1 - \prod_i (1 - P(\text{collision} \mid \xi(t), T_i(t)))$ -- cumulative collision probability,
  \item $R(\xi(t), t)$ -- predicted environmental risk field,
  \item $U(\xi(t)) = \sum_{i \in \mathcal{N}(\xi(t))} u_i \exp\left(-\frac{\|\xi(t) - \mathbf{p}_i\|^2}{2\sigma^2}\right)$ -- map uncertainty penalty,
  \item $\lambda_H, \lambda_R, \lambda_U$ -- weighting hyperparameters.
\end{itemize}

The selected action $\mathbf{a}^*$ minimizes the evaluated path cost:
\begin{equation}
  \mathbf{a}^* = \arg\min_{\mathbf{a} \in \mathcal{A}} J(\xi(\mathbf{a}))
\end{equation}
over admissible maneuvers (altitude change, heading adjustment, hovering, backtracking, or active information gathering)~\cite{pomianek2026s4dtam, cai2024neusis, yue2025airground, tordesillas2022faster}.

\subsection{Active Perception}
\label{subsec:active_perception}

The active perception module selects sensory action $\mathbf{a}_{\text{sense}}^*$ maximizing uncertainty reduction per unit cost:
\begin{equation}
  \mathbf{a}_{\text{sense}}^* = \arg\max_{\mathbf{a} \in \mathcal{A}_{\text{sense}}} \frac{I(T; \mathbf{z} \mid \mathbf{a})}{C(\mathbf{a})}
\end{equation}
where mutual information is approximated by token entropy reduction:
\begin{equation}
  I(T; \mathbf{z} \mid \mathbf{a}) \approx H(T) - \mathbb{E}_{\mathbf{z} \mid \mathbf{a}}[H(T \mid \mathbf{z})]
\end{equation}
and $C(\mathbf{a})$ denotes computational and gimbal energy costs~\cite{pomianek2026s4dtam, tao2024active, tian2024same}.

\subsection{Metrological Error Analysis of the S4D-TAM System}
\label{subsec:metrological_analysis}

\subsubsection{Sensor Layer Error Model}
Each sensor measurement $\mathbf{z}_i(t)$ conforms to:
\begin{equation}
  \mathbf{z}_i(t) = h_i(\mathbf{x}(t)) + \mathbf{n}_i(t) + \mathbf{b}_i(t)
\end{equation}
with zero-mean noise $\mathbf{n}_i(t) \sim \mathcal{N}(\mathbf{0}, \mathbf{R}_i)$ and time-varying bias $\mathbf{b}_i(t)$.

For IMU sensors (gyroscope and accelerometer):
\begin{align}
  \boldsymbol{\omega}(t) &= \boldsymbol{\omega}_{\text{true}}(t) + \mathbf{b}_g(t) + \mathbf{n}_g(t), \quad \mathbf{n}_g \sim \mathcal{N}(\mathbf{0}, \sigma_g^2 \mathbf{I}_3) \\
  \mathbf{a}(t) &= \mathbf{a}_{\text{true}}(t) + \mathbf{b}_a(t) + \mathbf{n}_a(t), \quad \mathbf{n}_a \sim \mathcal{N}(\mathbf{0}, \sigma_a^2 \mathbf{I}_3)
\end{align}
Biases follow Brownian motion random walks: $\dot{\mathbf{b}}_g = \mathbf{n}_{bg}$ and $\dot{\mathbf{b}}_a = \mathbf{n}_{ba}$. IMU dead-reckoning drift grows cubically over unconstrained intervals:
\begin{equation}
  \sigma_{\text{IMU}}(t) = \frac{1}{2}\sigma_a t^2 + \frac{1}{6}\sigma_{ba} t^3
\end{equation}

In degraded GNSS regimes, signal quality $q_{\text{GNSS}}(t) \in [0, 1]$ scales measurement covariance:
\begin{equation}
  \mathbf{R}_{\text{GNSS}}(t) = \frac{\mathbf{R}_{\text{GNSS}}^{(0)}}{q_{\text{GNSS}}(t) + \varepsilon}
\end{equation}
Visual odometry drift accumulates with traveled distance $d$:
\begin{equation}
  \sigma_{\text{VO}}(d) = \alpha_{\text{VO}} d + \beta_{\text{VO}} \sqrt{d}
\end{equation}
Barometric altimeters introduce slow thermal bias drift:
\begin{equation}
  z_{\text{baro}}(t) = z_{\text{true}}(t) + b_{\text{baro}} + n_{\text{baro}}(t)
\end{equation}

\subsubsection{Error Propagation in the Factor Graph}
Let $\boldsymbol{\xi} = [\mathbf{p}^\top, \mathbf{v}^\top, \boldsymbol{\phi}^\top, \mathbf{b}_g^\top, \mathbf{b}_a^\top]^\top \in \mathbb{R}^{15}$ denote the state vector. The Cram\'er-Rao Lower Bound (CRLB) bounds estimator covariance:
\begin{equation}
  \mathrm{Cov}(\boldsymbol{\xi}) \ge \mathbf{J}^{-1}(\boldsymbol{\xi})
\end{equation}
where $\mathbf{J}(\boldsymbol{\xi}) = \sum_i \mathbf{H}_i^\top \mathbf{R}_i^{-1} \mathbf{H}_i + \mathbf{J}_{\text{prior}}$ is the Fisher Information Matrix (FIM). Marginal position variance satisfies:
\begin{equation}
  \sigma_p^2 \ge [\mathbf{J}^{-1}(\boldsymbol{\xi})]_{pp}
\end{equation}

Fusing all factors yields cumulative information:
\begin{equation}
  \mathbf{J}_{\text{fused}} = \mathbf{J}_{\text{IMU}} + \mathbf{J}_{\text{VO}} + q_{\text{GNSS}}\mathbf{J}_{\text{GNSS}} + \mathbf{J}_{\text{alt}} + \mathbf{J}_{\text{sem}} + \mathbf{J}_{\text{ref}}
\end{equation}

\subsubsection{Quantization and Discretization Error of Tokens}
Spatial quantization into grid cells of cell length $\ell$ introduces bounded centroid error $\boldsymbol{\varepsilon}_q = \mathbf{p} - \mathbf{c}_k$ with $\|\boldsymbol{\varepsilon}_q\| \le \frac{\ell\sqrt{3}}{2}$. Across three dimensions:
\begin{equation}
  \sigma_{q,1\text{D}}^2 = \frac{\ell^2}{12}, \qquad \sigma_q^2 = \frac{\ell^2}{4}
\end{equation}
Semantic assignment error is measured by the conditional entropy:
\begin{equation}
  \varepsilon_{\text{sem}}^{(k)} = -\sum_{c=1}^C p(c \mid \mathcal{T}_k) \log p(c \mid \mathcal{T}_k) = H(c \mid \mathcal{T}_k)
\end{equation}

\subsubsection{Lower Bound on Localization Error}
Under complete GNSS denial ($q_{\text{GNSS}} = 0$), position error is bounded by:
\begin{equation}
  \sigma_p^2(t) = \sigma_{p,0}^2 + \alpha_{\text{VO}}^2 d(t)^2 + \gamma_{\text{sem}} N_{\text{loop}}^{-1}
\end{equation}
where $N_{\text{loop}}$ is the number of semantic loop closures and $\gamma_{\text{sem}}$ is the semantic correction factor.

\subsubsection{Occupancy Forecasting Error}
Modeling token propagation as a stochastic Gauss-Markov process:
\begin{equation}
  \mathbf{\Sigma}_{\text{pred}}(\Delta t) = \mathbf{\Phi}(\Delta t) \mathbf{\Sigma}_k(t) \mathbf{\Phi}(\Delta t)^\top + \mathbf{Q}_{\text{dyn}}(\Delta t)
\end{equation}
Binary occupancy prediction accuracy is evaluated via the Brier Score:
\begin{equation}
  \mathrm{BS}(\Delta t) = \frac{1}{|\mathcal{V}|} \sum_{\mathbf{x} \in \mathcal{V}} \left( \mathcal{O}(\mathbf{x}, t+\Delta t) - \mathcal{O}_{\text{true}}(\mathbf{x}, t+\Delta t) \right)^2
\end{equation}

\subsubsection{Cumulative Error Budget}
Table~\ref{tab:error_budget} presents the cumulative error budget for the target UAV platform operating in GNSS-degraded conditions.

\begin{table}[htbp]
  \centering
  \caption{Cumulative error budget of the S4D-TAM system layers.}
  \label{tab:error_budget}
  \begin{tabular}{llll}
    \toprule
    \textbf{Error Source} & \textbf{Characteristic} & \textbf{Typical Value ($1\sigma$)} & \textbf{Time/Distance Dependency} \\
    \midrule
    IMU noise (accelerometer) & White noise $\sigma_a$ & $0.05\,\text{m/s}^2$ & $\sigma \propto t^2$ \\
    IMU bias drift & Random walk $\sigma_{ba}$ & $0.002\,\text{m/s}^2/\sqrt{\text{s}}$ & $\sigma \propto t^{3/2}$ \\
    Visual odometry & Linear drift $\alpha_{\text{VO}}$ & $1\%$ per meter & $\sigma \propto d$ \\
    GNSS (full signal) & $\sigma_{\text{GNSS}}^{(0)}$ & $1.5\,\text{m}$ (horizontal) & Constant \\
    GNSS (degraded) & $\sigma_{\text{GNSS}}(q)$ & $5$--$50\,\text{m}$ & $\propto 1/q_{\text{GNSS}}$ \\
    Barometric altimeter & $\sigma_{\text{baro}} + \text{bias}$ & $0.3\,\text{m} + 2\,\text{m}$ & Thermal drift \\
    Token quantization & $\sigma_q$ & $\ell/2$ ($0.5\,\text{m}$ for $\ell=1\,\text{m}$) & Constant \\
    Semantic classification & Classifier entropy $H(c \mid \mathcal{T})$ & $0.1$--$0.8\,\text{bit}$ & Bounded \\
    Occupancy prediction ($\Delta t = 1\,\text{s}$) & $\mathrm{BS}(\Delta t)$ & $0.05$--$0.15$ & $\propto \Delta t$ \\
    Total error (IMU+VO only, $d=100\,\text{m}$) & CRLB $\sigma_p$ & $\approx 2$--$5\,\text{m}$ & $\propto d$ \\
    Total error (IMU+VO+sem, $N_{\text{loop}}=5$) & CRLB $\sigma_p$ & $\approx 0.5$--$1.5\,\text{m}$ & Loop bounded \\
    \bottomrule
  \end{tabular}
\end{table}

\subsection{Extended Semantic Token Model}
\label{subsec:extended_token}

The generalized semantic-temporal token is defined as a 9-tuple:
\begin{equation}
  \mathcal{T}_k = \langle \boldsymbol{\mu}_k, \mathbf{\Sigma}_k, \mathbf{p}_c^{(k)}, \mathcal{O}_k, \delta_k, \mathbf{v}_k, \mathcal{R}_k, \Omega_k, \rho_k \rangle
\end{equation}
where:
\begin{itemize}
  \item $\boldsymbol{\mu}_k \in \mathbb{R}^3$ -- estimated centroid in world coordinates,
  \item $\mathbf{\Sigma}_k \in \mathbb{R}^{3 \times 3}$ -- spatial covariance matrix,
  \item $\mathbf{p}_c^{(k)} = (p(c_j \mid \mathcal{T}_k))_{j=1}^C$ -- categorical semantic class distribution,
  \item $\mathcal{O}_k = \{(t_i, s_i, w_i, \mathbf{z}_i)\}_{i=1}^{n_k}$ -- observation history with timestamps, sensor IDs, weights, and measurements,
  \item $\delta_k \in \{\text{static}, \text{dynamic}, \text{periodic}\}$ -- motion state classification,
  \item $\mathbf{v}_k \in \mathbb{R}^3$ -- linear velocity vector,
  \item $\mathcal{R}_k = \{r_{kl}\}$ -- spatial-semantic relations to neighboring tokens,
  \item $\Omega_k$ -- geometric occlusion cone,
  \item $\rho_k \in [0, 1]$ -- flight safety significance coefficient.
\end{itemize}

Observations in $\mathcal{O}_k$ are weighted according to sensor characteristics, detection confidence, and temporal decay:
\begin{equation}
  w_i = w_{\text{base}}(s_i) \cdot \mathrm{conf}_i \cdot \exp(-\lambda_s (t - t_i))
\end{equation}
Centroid and covariance updates follow weighted information fusion:
\begin{equation}
  \mathbf{\Sigma}_k^{-1} = \sum_{i=1}^{n_k} w_i \mathbf{\Sigma}_i^{-1}, \qquad \boldsymbol{\mu}_k = \mathbf{\Sigma}_k \sum_{i=1}^{n_k} w_i \mathbf{\Sigma}_i^{-1} \mathbf{z}_i^{(k)}
\end{equation}

Dynamic classification tests displacement sequence $\Delta \boldsymbol{\mu}_k^{(i)} = \boldsymbol{\mu}_k(t_i) - \boldsymbol{\mu}_k(t_{i-1})$:
\begin{equation}
  \delta_k = \begin{cases}
    \text{static}, & \text{if } \|\bar{\Delta \boldsymbol{\mu}}_k\| < \varepsilon_s \text{ and } \sigma_{\Delta} < \varepsilon_\sigma \\
    \text{periodic}, & \text{if } \mathrm{FFT}(\|\Delta \boldsymbol{\mu}_k^{(i)}\|) \text{ has dominant frequency } f^* \\
    \text{dynamic}, & \text{otherwise}
  \end{cases}
\end{equation}

Spatial relation $r_{kl} = \langle d_{kl}, \mathbf{e}_{kl}, \psi_{kl}, \alpha_{kl} \rangle$ defines edge attributes in scene graph $\mathcal{G} = (\mathcal{V}, \mathcal{E})$. Occlusion cone $\Omega_k(\mathbf{x}_{\text{UAV}})$ models unobserved volumes behind obstacles:
\begin{equation}
  \Omega_k(\mathbf{x}_{\text{UAV}}) = \left\{ \mathbf{x} \in \mathbb{R}^3 : \exists \lambda > 1, \mathbf{x} = \mathbf{x}_{\text{UAV}} + \lambda(\boldsymbol{\mu}_k - \mathbf{x}_{\text{UAV}}), \|\mathbf{x} - \mathbf{axis}_k\| < r_k \right\}
\end{equation}

Safety significance is evaluated as:
\begin{equation}
  \rho_k = \rho_{\text{class}}(c_k^*) \cdot (1 - \exp(-\beta_\rho n_k)) \cdot f_{\text{context}}(\mathcal{T}_k, \mathbf{x}_{\text{UAV}})
\end{equation}

\subsection{Bayesian Inference about Future Scene States}
\label{subsec:bayesian_inference}

\emph{Implementation status.} The current executable reference uses a strictly causal probabilistic occupancy and motion forecaster driven only by observations available up to the prediction time. It produces Bernoulli occupancy probabilities, flow means, uncertainty estimates, observability masks, and physically timed forecast horizons. The Dynamic Bayesian Network (DBN) and particle approximation below define the target probabilistic extension for multimodal future-state uncertainty.

The scene state at time $t$ is modeled as $S_t = \{\mathcal{T}_k(t)\}_{k=1}^K$ with transition factorization:
\begin{equation}
  p(S_{t+\Delta t} \mid S_t) = \prod_{k=1}^K p(\mathcal{T}_k(t+\Delta t) \mid \mathcal{T}_k(t), \delta_k, \mathcal{G}_t)
\end{equation}
Individual state transitions follow:
\begin{equation}
  p(\boldsymbol{\mu}_k(t+\Delta t) \mid \boldsymbol{\mu}_k(t), \delta_k) = \begin{cases}
    \mathcal{N}(\boldsymbol{\mu}_k, \mathbf{\Sigma}_k + \mathbf{Q}_s \Delta t), & \delta_k = \text{static} \\
    \mathcal{N}(\boldsymbol{\mu}_k + \mathbf{v}_k \Delta t, \mathbf{\Sigma}_k + \mathbf{Q}_d \Delta t), & \delta_k = \text{dynamic} \\
    \sum_j \alpha_j \mathcal{N}(\boldsymbol{\mu}_k^{(j)}, \mathbf{\Sigma}_k^{(j)}), & \delta_k = \text{periodic}
  \end{cases}
\end{equation}

Particle filtering approximates the state posterior $p(S_t \mid \mathbf{z}_{1:t}) \approx \sum_{i=1}^{N_p} w_t^{(i)} \delta(S_t - \mathbf{s}_t^{(i)})$ with effective sample size resampling:
\begin{equation}
  N_{\text{eff}} = \frac{1}{\sum_{i=1}^{N_p} (w_t^{(i)})^2} < \frac{N_p}{2}
\end{equation}

Occluded zone occupancy is inferred via:
\begin{equation}
  p(\text{occupied}(\mathbf{x}) \mid \mathcal{T}_k) = p_0(c_k^*) + (1 - p_0(c_k^*)) \exp\left(-\frac{d(\mathbf{x}, \partial \Omega_k)^2}{2\sigma_\Omega^2}\right), \quad \mathbf{x} \in \Omega_k
\end{equation}
Anomalies trigger alert modes when Mahalanobis distance exceeds threshold:
\begin{equation}
  D_M(t) = \sqrt{(S_t^{\text{obs}} - S_t)^\top \mathbf{\Sigma}_S^{-1} (S_t^{\text{obs}} - S_t)} > D_{\text{alarm}}
\end{equation}

\subsection{Bayesian Risk Estimation and Trajectory Planning}
\label{subsec:bayesian_risk_planning}

\emph{Implementation status.} The executable reference planner currently uses deterministic finite-horizon beam search over admissible acceleration controls with explicit kinematic limits and decomposed collision-risk, energy, time, goal-progress, and information-value terms. MPPI remains a planned backend for richer dynamics and will be introduced behind the same planner contract.

The risk field $\mathcal{R}(\mathbf{x}, t)$ accumulates risk across tokens:
\begin{equation}
  \mathcal{R}(\mathbf{x}, t) = \sum_{k=1}^K \rho_k \cdot p_{\text{occ}}(\mathbf{x}, \mathcal{T}_k, t) \cdot p_{\text{coll}}(\mathbf{x}, \mathcal{T}_k)
\end{equation}
with collision probability:
\begin{equation}
  p_{\text{coll}}(\mathbf{x}, \mathcal{T}_k) = 1 - \Phi\left(\frac{d(\mathbf{x}, \boldsymbol{\mu}_k) - r_{\text{safe}}}{\sqrt{\sigma_{k,r}^2 + \sigma_{\text{UAV}}^2}}\right)
\end{equation}
and occlusion augmentation $\mathcal{R}(\mathbf{x}, t) \leftarrow \mathcal{R}(\mathbf{x}, t) + \kappa_\Omega \cdot \mathbf{1}_{\mathbf{x} \in \bigcup \Omega_k} \cdot \max_k p_0(c_k^*)$.

The multi-objective trajectory cost is:
\begin{equation}
  J(\tau) = w_L L(\tau) + w_E E(\tau) + w_R \int_\tau \mathcal{R}(\mathbf{x}(t), t) dt + w_U \int_\tau \mathcal{U}(\mathbf{x}(t), t) dt - w_I \mathcal{I}(\tau)
\end{equation}
where information gain $\mathcal{I}(\tau) = H(S_t) - \mathbb{E}_\tau[H(S_t \mid \mathbf{z}_{1:\tau})]$.

In the target MPPI formulation:
\begin{equation}
  \tau^* \approx \frac{\sum_{m=1}^M \exp\left(-\frac{1}{\lambda_{\text{MPPI}}} J(\tau_m)\right) \tau_m}{\sum_{m=1}^M \exp\left(-\frac{1}{\lambda_{\text{MPPI}}} J(\tau_m)\right)}
\end{equation}

Operational risk defines four UAV flight modes:
\begin{equation}
  \text{UAV Mode} = \begin{cases}
    \text{nominal}, & \mathcal{R}(\mathbf{x}_{\text{UAV}}, t) < \mathcal{R}_{\text{safe}} \\
    \text{alert}, & \mathcal{R}_{\text{safe}} \le \mathcal{R} < \mathcal{R}_{\text{alert}} \\
    \text{avoidance}, & \mathcal{R}_{\text{alert}} \le \mathcal{R} < \mathcal{R}_{\text{avoid}} \\
    \text{emergency}, & \mathcal{R} \ge \mathcal{R}_{\text{avoid}}
  \end{cases}
\end{equation}

\subsection{Adaptive World Model with Hierarchical Resolution}
\label{subsec:adaptive_world_model}

\emph{Implementation status.} The current token memory enforces bounded history and configurable limits on token count, map memory, and update time. The full adaptive octree and Level-of-Detail (LOD) storage policy described below represents a planned world-model backend.

Operational space is structured as an octree $\mathcal{O}$ with hierarchical voxel resolutions $\ell \in \{\ell_1 = 0.1\,\text{m}, \ell_2 = 0.5\,\text{m}, \ell_3 = 2\,\text{m}, \ell_4 = 10\,\text{m}\}$ mapped to LOD levels:
\begin{itemize}
  \item \textbf{LOD 0}: $(\boldsymbol{\mu}_k, \rho_k, c_k^*)$ -- centroid, risk, dominant class ($64\,\text{B}$).
  \item \textbf{LOD 1}: LOD 0 + subsampled points ($\le 64$ pts) + $\mathbf{p}_c^{(k)}$ ($512\,\text{B}$).
  \item \textbf{LOD 2}: LOD 1 + attention embedding + observation history $\mathcal{O}_k$ ($4\,\text{KB}$).
  \item \textbf{LOD 3}: LOD 2 + dynamic prediction model + relations $\mathcal{R}_k$ + occlusion cone $\Omega_k$ ($32\,\text{KB}$).
\end{itemize}

Optimal token resolution adapts to distance, risk, and CPU availability:
\begin{equation}
  \ell_k = \ell_1 \cdot 2^{L_k}, \qquad L_k = \left\lfloor L_{\max} \cdot f_{\text{dist}}(d_k) \cdot f_{\text{risk}}(\rho_k)^{-1} \cdot f_{\text{comp}}(C_{\text{avail}})^{-1} \right\rfloor
\end{equation}
Token update priority is scheduled by:
\begin{equation}
  \pi_k = \alpha_\pi \rho_k + \beta_\pi \frac{1}{\lambda_{\min}(\mathbf{\Sigma}_k) + \varepsilon} + \gamma_\pi (1 - \exp(-\nu(t - t_{\text{last},k}))) + \delta_\pi \mathbf{1}_{\delta_k \ne \text{static}}
\end{equation}

Computational budget allocation enforces cascade degradation under CPU load:
\begin{equation}
  C_{\text{total}} = C_{\text{sensor}} + C_{\text{SLAM}} + C_{\text{token}} + C_{\text{attn}} + C_{\text{pred}} + C_{\text{plan}}
\end{equation}

\subsection{Computational Resource Management and Bottleneck Elimination}
\label{subsec:resource_management}

\emph{Implementation status.} Structured telemetry and resource budgets expose update-time, token-count, memory, lifecycle, and decision events needed for profiling and replay. Predictive queue-overflow control, EDF-RT scheduling, and survival-mode control define target extensions.

Pipeline throughput is constrained by the bottleneck stage:
\begin{equation}
  \Theta_{\text{sys}} = \frac{1}{\max_i T_i}, \qquad T_{\text{cycle}} = \sum_{i=1}^6 T_i + T_{\text{comm}}
\end{equation}

Queue overflow probability is modeled as a Poisson arrival process:
\begin{equation}
  P(q_i(t+\Delta t) > q_{\max}) = 1 - \sum_{n=0}^{q_{\max}-q_i(t)} \frac{(\lambda_i \Delta t)^n}{n!} e^{-\lambda_i \Delta t}
\end{equation}

Dynamic processing frequency adapts as $f_{\text{proc}}(t) = \min(f_{\text{nom}}, \frac{C_{\text{avail}}(t)}{C_{\text{per}}})$. Memory is governed by TTL deactivation, spatial-semantic merging ($\|\boldsymbol{\mu}_k - \boldsymbol{\mu}_l\| < \ell_k$ and $D_{\text{KL}}(\mathbf{p}_c^{(k)} \| \mathbf{p}_c^{(l)}) < \varepsilon_{\text{KL}}$), and LOD compression.

Real-time scheduling prioritizes safety-critical tasks via:
\begin{equation}
  \mathrm{Priority}_i(t) = \frac{\rho_i^{\max}}{T_{\text{deadline},i}(t) - t} + \kappa_{\text{safety}} \cdot \mathbf{1}_{i \in \mathcal{C}_{\text{safety}}}
\end{equation}
guaranteeing minimal CPU reservations for sensor acquisition, SLAM tracking, and collision avoidance.

\section{Research Hypotheses}
\label{sec:hypotheses}

The confirmatory mechanism study uses seven preregistered hypotheses. Each hypothesis compares the full S4D-TAM model with a variant that disables exactly one mechanism while data, seeds, preprocessing, optimization budget, and all other settings remain fixed. External systems are deliberately excluded from this causal ablation family.

\begin{itemize}
  \item \textbf{H1: Semantics improve navigation in scenarios containing dynamic objects.} The primary endpoint is mission success; collisions per kilometer and semantic mIoU are secondary outcomes.
  \item \textbf{H2: Temporal state improves 1\,s future-state prediction.} The primary endpoint is forecast occupancy IoU at 1\,s; flow EPE and temporal flip rate are secondary outcomes.
  \item \textbf{H3: Calibrated uncertainty reduces navigation risk.} The primary endpoint is collisions per kilometer; ECE, NLL, and near misses are secondary outcomes.
  \item \textbf{H4: Topological support increases mission success.} Path efficiency and ATE RMSE are secondary outcomes.
  \item \textbf{H5: Reference-map support reduces localization error.} The primary endpoint is ATE RMSE; RPE and final drift percentage are secondary outcomes.
  \item \textbf{H6: Predictive risk modeling reduces collisions.} The primary endpoint is collisions per kilometer; near misses and minimum clearance are secondary outcomes.
  \item \textbf{H7: Token lifecycle management reduces computational cost without materially degrading navigation quality.} Confirmation requires both lower p95 latency and non-inferior mission success with a preregistered margin of 2 percentage points; map bytes and peak RSS are secondary outcomes.
\end{itemize}

Hypotheses H1--H6 are superiority hypotheses in the preregistered direction. H7 combines efficiency superiority with mission-success non-inferiority. The seven confirmatory decisions form one family and use the Holm procedure defined in Section~\ref{subsec:statistical_analysis}.

\section{Proposed Experimental Validation}
\label{sec:experiments}

\subsection{Datasets}
\label{subsec:datasets}

The validation framework targets four distinct benchmark datasets: TartanAir~\cite{wang2020tartanair}, Blackbird, MARSIM~\cite{kong2023marsim}, and AeroVerse. Each dataset is converted to the common \texttt{SequenceData} contract with immutable source version, license record, coordinate convention, calibration metadata, and SHA-256 manifest. The cohort identifier \texttt{frozen-2026-01} denotes the study cohort rather than a vendor release and resolves to exact source versions and file hashes.

Data are split at the environment or flight level, never across adjacent temporal frames. The training split is the only source of learned weight updates; calibration splits are strictly reserved for hyperparameter thresholds, normalization, early stopping, and uncertainty calibration; the test split remains inaccessible until code, artifacts, and analysis scripts are frozen.

The preregistered offline mechanism study uses 20 preselected test sequences per dataset, yielding 80 paired sequence-level units. Systems validation additionally defines 30 paired Software-in-the-Loop (SIL) scenarios with five random seeds, 20 Hardware-in-the-Loop (HIL) scenarios with five seeds, and 20 paired controlled real-flight missions per variant when safety gates remain open.

A dedicated UAV dataset remains desirable for synchronized RGB, thermal, IMU, GNSS, point cloud, elevation/reference maps, static/dynamic ground truth labels, visibility masks, and controlled sensor degradation scenarios. Metrics that require unavailable ground truth are reported as unavailable rather than silently omitted or imputed.

\subsection{Two-Level Comparison Design}
\label{subsec:comparison_design}

The evaluation is decoupled into two non-interchangeable levels. Both levels use common dataset and metric contracts where applicable, but they address distinct scientific questions and produce separate statistical reports.

\subsubsection{External System Comparison}
\label{subsubsec:external_comparison}

The external comparison evaluates whether the complete S4D-TAM system is competitive with independently implemented navigation and SLAM systems under identical data, hardware constraints, sensor availability, and evaluator definitions.

The mandatory primary baseline set comprises:
\begin{itemize}
  \item \textbf{ORB-SLAM3},
  \item \textbf{VINS-Mono}~\cite{qin2018vinsmono},
  \item \textbf{FAST-LIO2}~\cite{xu2022fastlio2},
  \item \textbf{LIO-SAM}~\cite{shan2020liosam}.
\end{itemize}
Each baseline is pinned to an upstream git commit revision, fixed containerized environment, calibration profile, sensor-topic mapping, loop-closure policy, warm-up policy, and hardware execution budget. Results enter the benchmark exclusively through the common \texttt{AlgorithmResult} schema.

This level measures complete-system competitiveness. Observed differences relative to external baselines cannot be interpreted as the isolated causal effect of an individual S4D-TAM component.

\subsubsection{Internal Mechanism Study}
\label{subsubsec:internal_study}

The internal comparison isolates mechanism contributions using the full system and seven single-component ablation variants:
\begin{itemize}
  \item \texttt{H1\_no\_semantics},
  \item \texttt{H2\_no\_temporal\_state},
  \item \texttt{H3\_no\_calibrated\_uncertainty},
  \item \texttt{H4\_no\_topology},
  \item \texttt{H5\_no\_reference\_map},
  \item \texttt{H6\_no\_risk\_prediction},
  \item \texttt{H7\_no\_token\_lifecycle}.
\end{itemize}
Each variant differs from the full system by exactly one switch and utilizes its own frozen trained artifact. External systems are excluded from this matrix to avoid confounding complete-system identity with the causal effect of removing an individual mechanism.

\subsection{Metrics}
\label{subsec:metrics}

Primary endpoints are preregistered:
\begin{itemize}
  \item \textbf{Localization}: Absolute Trajectory Error (ATE) RMSE after $\mathrm{SE}(3)$ rigid alignment without scale optimization.
  \item \textbf{Forecasting}: Occupancy IoU at $1\,\text{s}$ (and $3\,\text{s}$ for multi-step characterization).
  \item \textbf{Safety}: Mission success rate and collision count per kilometer.
  \item \textbf{Efficiency}: 95th percentile (p95) processing latency and persistent map-memory footprint.
\end{itemize}

Secondary localization metrics include Relative Pose Error (RPE) translation RMSE, final drift in meters and percentage of traveled distance, rotation RPE, relocalization success rate, and time-to-relocalize. Semantic metrics encompass mIoU and macro F1 score. Forecasting secondary metrics include flow End-Point Error (EPE), Brier score, Negative Log-Likelihood (NLL), and Expected Calibration Error (ECE). Navigation metrics include near misses, minimum obstacle clearance, and path efficiency.

\subsection{Confirmatory Protocol and Statistical Analysis}
\label{subsec:statistical_analysis}

The preregistered H1--H7 mechanism study conducts sequence- or mission-level inference rather than treating individual frames as independent samples. Continuous outcomes are evaluated via paired mean differences and $95\%$ BCa bootstrap confidence intervals using $10{,}000$ resamples. Mission success is modeled as a paired scenario-level binary probability outcome, while distance-normalized event counts utilize Poisson or negative-binomial generalized linear models with log-distance offset.

The seven H1--H7 confirmatory decisions form one family controlled by the sequential Holm step-down procedure at family-wise error rate $\alpha = 0.05$. Hypotheses H1--H6 require the adjusted $p$-value and confidence interval to demonstrate superiority in the preregistered direction. Hypothesis H7 additionally requires non-inferiority in mission success within a $2$ percentage-point margin while demonstrating p95 latency reduction.

\section{Discussion}
\label{sec:discussion}

The proposed S4D-TAM system represents an attempt to transition from a classical geometric map to a predictive spatial memory for autonomous UAVs. The primary advantage is \emph{amortized perception} -- the computational cost of deep image analysis is incurred once during token instantiation and update, while the synthesized semantic representation remains persistently accessible in spatial memory across subsequent planning iterations~\cite{pomianek2026s4dtam}. This structure minimizes computational redundancy where identical raw frames are repeatedly processed by independent downstream modules~\cite{rosen2024scene, ding2025survey}.

The hierarchical token representation enables storing expansive, regular terrain surfaces at coarse resolutions while concentrating high-resolution representation on structural edges, obstacles, and dynamic entities~\cite{li2024hislam, tian2024same, xiang2021customizable, zhang2021you, he2024efficient}. This paradigm aligns with recent advances in superpoint and octree graph transformers showing that aggregating geometric points into structured primitives substantially reduces model parameter footprint and attention complexity.

The spatiotemporal attention mechanism models long-range dependencies across the environment while constraining computational overhead by restricting query-key interaction to local topological neighborhoods and globally salient landmarks~\cite{yu2025driveoccworld, silva2024s2tpvformer, liao2025i2world, pcpnet2024}. Uncertainty-gated attention prevents noisy or degraded measurements from destabilizing the persistent world state, which is crucial under intermittent sensing or partial visibility~\cite{pomianek2026s4dtam, cai2024neusis}.

Occupancy and risk forecasting in unobserved regions marks a transition from reactive detection to proactive prediction~\cite{pomianek2026s4dtam, yu2025driveoccworld, mei2025irwm, liao2025i2world, zheng2024gaussianad, liu2025hybrid}. Rather than hallucinating unseen obstacles, the system computes rigorous probabilistic hazard distributions informed by topological continuity, semantic class priors, and historical dynamic flows. This enables defensive trajectory planning that anticipates emerging threats before they enter direct line of sight.

Reference map co-registration provides robust drift mitigation during prolonged GNSS dropouts~\cite{liu2022largescale, yue2025airground, radwan2024uav}. Map residuals simultaneously serve as diagnostic indicators for localization errors, structural scene modifications, and GNSS spoofing anomalies~\cite{pomianek2026s4dtam}.

Active perception dynamically allocates sensor and computing bandwidth by maximizing information gain relative to energetic cost~\cite{pomianek2026s4dtam, tao2024active, tian2024same}. Modality selection (e.g., triggering thermal sensors in low illumination or throttling RGB deep encoders) ensures predictable real-time execution on resource-constrained embedded platforms.

\subsection{Limitations}
\label{subsec:limitations}

Current limitations of the architecture include:
\begin{enumerate}
  \item \textbf{System Complexity}: Tight integration and synchronization across SLAM, open-vocabulary segmentation, transformer backbones, forecasting models, and trajectory planners require strict lifecycle management.
  \item \textbf{Computational Demand}: Although perception is amortized, transformer attention and dense probabilistic forecasting remain computationally demanding on embedded hardware.
  \item \textbf{Reference Map Dependency}: The efficacy of reference alignment depends on the accuracy and resolution of prior elevation or point cloud databases.
  \item \textbf{Uncertainty Calibration}: Ensuring reliable calibration of token confidence and risk bounds across diverse environmental conditions requires extensive calibration validation.
  \item \textbf{Domain Generalization}: Adapting semantic priors across extreme shifts in weather, vegetation, and urban topography may require targeted domain adaptation.
\end{enumerate}

\section{Conclusions and Future Research Directions}
\label{sec:conclusions}

This paper has presented S4D-TAM (Semantic 4D Token Attention Map), a novel autonomous navigation architecture for unmanned aerial vehicles operating in GNSS-degraded environments. The foundational contribution is the persistent 4D semantic token map, unifying localization, semantic scene understanding, spatiotemporal occupancy forecasting, risk estimation, and motion planning into a single coherent spatial representation.

The key conceptual contributions include:
\begin{enumerate}
  \item \textbf{4D Semantic Token Model}: Compactly integrating spatial geometry, continuous semantics, dynamic state, observation history, and calibrated uncertainty.
  \item \textbf{Hierarchical Multi-Layer Architecture}: Unifying multi-sensor fusion, SLAM front-end, tokenization, transformer attention, memory persistence, reference map alignment, probabilistic forecasting, risk-aware trajectory planning, and active perception.
  \item \textbf{Spatiotemporal Attention Mechanism}: Gating attention by spatial proximity, temporal decay, and uncertainty bounds to guarantee robust state estimation.
  \item \textbf{Proactive Occlusion and Risk Forecasting}: Inferring hazard probabilities in unobserved spatial volumes conditioned on semantic topology and dynamic flows.
  \item \textbf{Multi-Objective Risk-Aware Planning}: Optimizing flight trajectories against safety risk, path length, energy expenditure, and spatial uncertainty.
  \item \textbf{Information-Theoretic Active Perception}: Dynamically allocating sensor modalities and processing rates based on mutual information gain and energy constraints.
\end{enumerate}

Seven preregistered mechanism hypotheses (H1--H7) have been defined to causally isolate the impact of individual architectural features from complete-system baseline comparisons against ORB-SLAM3, VINS-Mono, FAST-LIO2, and LIO-SAM.

Future research directions will concentrate on:
\begin{enumerate}
  \item Executing the prospective two-level validation study across TartanAir, Blackbird, MARSIM, and AeroVerse datasets, progressing from offline benchmarks to SIL, HIL, and field flight tests.
  \item Embedded model optimization through post-training quantization, structured token pruning, and hardware-accelerated attention backends.
  \item Expanding multi-sensor modalities to millimeter-wave radar and ultrasound, and evaluating multi-UAV cooperative mapping teams.
  \item Integrating multimodal large language models (LLMs) for high-level semantic reasoning and natural language mission command.
\end{enumerate}

\section*{Acknowledgements}
\label{sec:acknowledgements}

The author thanks the reviewers for their valuable comments and the Department of Computer and Control Engineering of Rzesz\'ow University of Technology for research support.


% --- Bibliography ---
\bibliographystyle{IEEEtran}
\bibliography{references}

\end{document}
